% Substep 4 documentation (Newton / MuJoCo-Warp integration)
%
% This file documents the implementation work done in substep 4 of the tire-model
% porting plan: integrating the previously validated tire math (substep 2) and
% disc-terrain collision (substep 3) into Newton + MuJoCo-Warp via mjcb_control.
%
\section{Substep 4: Newton / MuJoCo-Warp Integration (Chrono-like Tires)}
\label{sec:substep4-mujoco-integration}

\subsection{Goal}
Integrate the Chrono-like tire pipeline into the Newton MuJoCo-Warp solver such that:
\begin{itemize}
  \item tire--terrain interaction is handled \emph{explicitly} (Chrono-style) via externally applied wrenches,
  \item the same core collision and tire-force logic used in production is validated by unit tests against Chrono (ground truth),
  \item the implementation supports \texttt{nworld=2} and is ready for warp-batched execution.
\end{itemize}
Chrono remains the ground-truth oracle; we do not modify Chrono to make tests pass.

\subsection{Design Choice: ``external tire forces'' via \texttt{mjcb\_control}}
\label{sec:substep4-design-choice}
MuJoCo computes contact forces as part of its constraint/optimization solve. Chrono's tire pipeline instead:
\begin{enumerate}
  \item queries terrain height/normal and computes disc--terrain contact geometry (contact patch, penetration depth),
  \item computes tire forces/moments as explicit functions of state and contact data,
  \item applies the resulting wrench at the wheel body.
\end{enumerate}

For massively-parallel RL, we implement the Chrono-style route by computing tire forces in a MuJoCo-Warp
\texttt{mjcb\_control} callback and writing them to \texttt{xfrc\_applied}:
\begin{itemize}
  \item This callback is invoked from MuJoCo(-Warp) \texttt{forward()}, and thus is evaluated \emph{per RK4 stage}
        (important for consistent integration).
  \item The wheel--terrain ``contact'' is not handled by MuJoCo's contact solver; instead, wheel contact is disabled
        (recommended: only for wheel geoms) and replaced by our explicit tire forces.
\end{itemize}

\subsection{What Was Added / Modified (Key Files)}
\begin{itemize}
  \item \texttt{newton/newton/\_src/vehicle/tires/mujoco\_tire\_module.py}:
    production tire module \texttt{MujocoFialaTireModule} that computes contact, tire forces, and writes \texttt{xfrc\_applied}.

  \item \texttt{newton/newton/\_src/vehicle/tires/disc\_terrain\_collision.py}:
    MuJoCo plane/hfield terrain queries and the warp-batched disc--terrain collision kernel used by the tire module.
    A core requirement addressed in the final review is that \emph{the same Warp collision implementation is used by production and tests}.

  \item \texttt{newton/newton/\_src/solvers/mujoco/solver\_mujoco.py}:
    installs the tire modules into \texttt{mujoco\_warp.mjcb\_control} once during initialization and
    disables wheel collisions (recommended Chrono-like usage).

  \item \texttt{newton/newton/tests/tires/test\_mujoco\_tire\_integration.py}:
    integration tests that compare MuJoCo-Warp + \texttt{MujocoFialaTireModule} results against Chrono ground truth.

  \item \texttt{newton/newton/tests/tires/test\_disc\_terrain\_collision.py}:
    unit tests comparing the production-path disc--terrain collision kernel against Chrono ground truth for MuJoCo plane/hfield terrain.
\end{itemize}

\subsection{Production Module: \texttt{MujocoFialaTireModule}}
\label{sec:substep4-mujoco-tire-module}
File: \texttt{newton/newton/\_src/vehicle/tires/mujoco\_tire\_module.py}

\subsubsection{Inputs and configuration}
The module is configured with:
\begin{itemize}
  \item a Chrono-like \texttt{FialaTire} instance loaded from a Chrono vehicle JSON file,
  \item a list of MuJoCo wheel body ids (\texttt{wheel\_body\_ids}),
  \item a MuJoCo terrain geom id (\texttt{terrain\_geom\_id}),
  \item a collision mode (\texttt{CollisionType}: single point / four points / envelope).
\end{itemize}

Two convenience constructors are provided:
\begin{itemize}
  \item \texttt{from\_json(...)}: provide ids directly.
  \item \texttt{from\_mujoco\_names(...)}: resolve body/geom names using MuJoCo name-to-id lookup and validate terrain geom types.
\end{itemize}
Only MuJoCo \texttt{plane} and \texttt{hfield} terrains are supported in the production query helper.

\subsubsection{Per-callback data flow}
On each callback (\texttt{apply(self, m, d)}), for each world and wheel:
\begin{enumerate}
  \item \textbf{Disc--terrain collision}:
    launch \texttt{\_mujoco\_disc\_terrain\_collision\_kernel} to compute:
    \texttt{in\_contact}, penetration \texttt{depth}, friction \texttt{mu},
    and the contact patch frame \texttt{(pos, x\_axis, y\_axis, z\_axis)}.

  \item \textbf{Contact kinematics}:
    compute wheel velocity components $(v_x, v_y, v_z)$ in the contact patch frame and the wheel spin $\omega$
    about the disc normal (Chrono wheel axis convention).

  \item \textbf{Tire forces/moments}:
    launch \texttt{fiala\_tire\_advance\_kernel} (from substep~2) to compute:
    $F_x, F_y, F_z, M_y, M_z$ and also the diagnostic slips $(\kappa, \alpha)$.

  \item \textbf{Apply wrench}:
    convert the tire-frame forces/moments into world-frame and write them into MuJoCo-Warp's
    \texttt{d.xfrc\_applied[world, body]}.
    The applied torque includes the moment shift from the contact patch to the wheel center:
    \[
      \boldsymbol{\tau}_{\text{center}} =
      \boldsymbol{\tau}_{\text{patch}} + (\mathbf{r} \times \mathbf{F})
    \]
    with $\mathbf{r}$ being the vector from wheel center to the effective application point used by Chrono.
\end{enumerate}

\paragraph{Wheel axis convention}
Chrono's wheel axis is the local $Y$ axis. In this integration, the disc normal is taken as the wheel $Y$ axis
in world coordinates (MuJoCo body rotation matrix column 1). This matches the convention used in the Chrono-like
collision pipeline (\texttt{disc\_normal}).

\subsection{Terrain Query Support (MuJoCo plane + hfield)}
\label{sec:substep4-terrain-query}
File: \texttt{newton/newton/\_src/vehicle/tires/disc\_terrain\_collision.py}

\subsubsection{Supported types}
The production terrain query helper \texttt{\_wp\_mujoco\_terrain\_height\_normal} supports:
\begin{itemize}
  \item \texttt{mjGEOM\_PLANE}: height and normal computed from geom pose and geom $z$-axis,
  \item \texttt{mjGEOM\_HFIELD}: height and normal computed using MuJoCo's grid layout + triangulation rule.
\end{itemize}
Unsupported types are treated as an error on the host side (validated in \texttt{from\_mujoco\_names} and at first \texttt{apply}).

\subsubsection{Heightfield mapping}
For a MuJoCo hfield:
\begin{itemize}
  \item local coordinates $(x,y)$ are mapped into grid cell indices $(r,c)$ using:
    \[
      u = \frac{x + s_x}{2s_x/(n_c-1)},\quad v = \frac{y + s_y}{2s_y/(n_r-1)}
    \]
  \item vertex heights are \texttt{hfield\_data[idx] * size\_z\_top},
  \item each cell is triangulated along the diagonal $v_{00}\rightarrow v_{11}$ (MuJoCo convention),
        and both height and normal are computed accordingly.
\end{itemize}
The helper supports translated and rotated hfields by transforming the query point into the hfield local frame
using the full 3D delta (not just the $x$/$y$ components).

\subsection{Deduplication Requirement: ``tests hit the production code path''}
\label{sec:substep4-dedup}
In the final review of substep~4, an important correctness requirement was enforced:
\begin{quote}
  The production kernels used by \texttt{MujocoFialaTireModule} must be the same kernels exercised in unit tests.
\end{quote}

\subsubsection{Single shared collision implementation}
Warp kernels cannot use Chrono-style OO polymorphism inside device code. The implemented solution is:
\begin{itemize}
  \item a single core Warp function \texttt{\_wp\_disc\_terrain\_collision(...)} implementing the Chrono
        1pt/4pt/envelope algorithms while querying terrain via \texttt{\_wp\_mujoco\_terrain\_height\_normal},
  \item a production kernel \texttt{\_mujoco\_disc\_terrain\_collision\_kernel} that computes disc center/normal
        from MuJoCo state and then calls the shared function,
  \item a test wrapper kernel \texttt{\_disc\_terrain\_collision\_mujoco\_kernel} that accepts externally provided
        disc center/normal arrays and calls the same shared function.
\end{itemize}
This matches Chrono's structure conceptually: a single collision algorithm calling a terrain query interface.

\subsection{Solver Integration: installing tire modules and disabling wheel contacts}
\label{sec:substep4-solver}
File: \texttt{newton/newton/\_src/solvers/mujoco/solver\_mujoco.py}

\subsubsection{Callback installation}
If tire modules are provided to the solver, the solver installs a chained MuJoCo-Warp control callback once
at initialization time:
\begin{itemize}
  \item any previous callback is called first (if present),
  \item then each tire module's \texttt{apply(m, d)} method is executed.
\end{itemize}
This avoids reassigning the callback every step and ensures failures are not hidden by callback restoration.

\subsubsection{Disabling MuJoCo contact handling for wheels}
For Chrono-style external tire forces, wheel geoms should not participate in MuJoCo contact.
The solver disables collisions for all geoms attached to the wheel bodies declared by the tire modules by setting:
\texttt{geom\_contype = 0} and \texttt{geom\_conaffinity = 0}.

\subsection{Tests}
\label{sec:substep4-tests}

\subsubsection{Disc--terrain collision tests (unit)}
File: \texttt{newton/newton/tests/tires/test\_disc\_terrain\_collision.py}
\begin{itemize}
  \item Launches the MuJoCo-path collision wrapper kernel (\texttt{\_disc\_terrain\_collision\_mujoco\_kernel})
        on CPU for \texttt{nworld=2}.
  \item Compares outputs (\texttt{in\_contact, depth, mu, contact frame}) against Chrono ground truth via the CLI oracle.
  \item Validates both MuJoCo plane and MuJoCo hfield (sinusoid is represented as an hfield, avoiding analytic test-only terrain in production code).
\end{itemize}

\subsubsection{MuJoCo-Warp + tire module integration tests}
File: \texttt{newton/newton/tests/tires/test\_mujoco\_tire\_integration.py}
\begin{itemize}
  \item \textbf{RK4 callback stage correctness}:
        verifies that \texttt{mjcb\_control} is called 4 times per step when using RK4 integration.
  \item \textbf{Plane ground truth comparison}:
        applies \texttt{MujocoFialaTireModule} and compares \texttt{xfrc\_applied} at the wheel to a Chrono-derived wrench.
        The test includes a contact and a no-contact world and asserts contact actually occurs in the contact case.
  \item \textbf{Heightfield ground truth comparison}:
        repeats the same idea on a MuJoCo hfield terrain. The ``sinusoid'' surface is encoded into the hfield.
        The test validates all collision modes.
\end{itemize}

\subsection{Running the Tests}
From the workspace root:
\begin{verbatim}
PYTHONPATH=./newton:./mujoco_warp \
  newton/.venv/bin/python -m unittest -v newton.tests.tires.test_disc_terrain_collision

PYTHONPATH=./newton:./mujoco_warp \
  newton/.venv/bin/python -m unittest -v newton.tests.tires.test_mujoco_tire_integration
\end{verbatim}

\paragraph{Warp cache path}
If Warp cannot write to \texttt{\$HOME/.cache/warp}, set:
\begin{verbatim}
export WARP_CACHE_PATH=$(pwd)/.warp_cache
\end{verbatim}
The tire tests also force a repo-local kernel cache at import time to make runs robust in sandboxed environments.

\paragraph{Chrono oracle binary}
The tests require the Chrono ground-truth CLI from substep~0.
If missing, build it as documented in:
\texttt{newton/newton/tests/tires/chrono\_gt/README.md}
or set \texttt{NEWTON\_CHRONO\_GT\_BIN} to an existing binary.

\subsection{Current Scope / Limitations}
\begin{itemize}
  \item Production terrain query supports only MuJoCo \texttt{plane} and \texttt{hfield} terrains.
  \item The current integration compares forces/moments at the wheel center; trajectory-level matching is deferred to substep~5.
  \item Additional Newton-facing ``wheel tagging'' and full vehicle authoring UX can be expanded as needed, but is not required for the unit-tested core pipeline.
\end{itemize}

